
\begin{figure}[htbp]
\centering
\resizebox{0.92\textwidth}{!}{%
\begin{tikzpicture}[font=\scriptsize, node distance=9mm]
  \tikzset{ap/.style={draw, rounded corners, fill=warnbg, align=center, minimum width=16mm, minimum height=7mm},
           sta/.style={draw, circle, fill=lightaccent, align=center, minimum size=8mm}}
  \draw[fill=black!3, draw=accent!60, rounded corners] (-2.1,-1.7) rectangle (2.1,1.7);
  \draw[fill=black!3, draw=accent!60, rounded corners] (4.0,-1.7) rectangle (8.2,1.7);
  \node at (0,1.45) {BSS 1};
  \node at (6.1,1.45) {BSS 2};

  \node[ap] (ap1) at (0,0) {AP1};
  \node[ap] (ap2) at (6.1,0) {AP2};
  \node[sta] (s1) at (-1.2,0.8) {STA};
  \node[sta] (s2) at (1.1,-0.9) {STA};
  \node[sta] (s3) at (4.9,-0.8) {STA};
  \node[sta] (s4) at (7.2,0.8) {STA};

  \draw[dashed] (s1) -- (ap1);
  \draw[dashed] (s2) -- (ap1);
  \draw[dashed] (s3) -- (ap2);
  \draw[dashed] (s4) -- (ap2);

  \draw[very thick, accent] (ap1) -- node[above]{DS -- Distribution System} (ap2);
  \draw[decorate, decoration={brace, amplitude=5pt}] (-2.1,-2.05) -- (8.2,-2.05) node[midway, below=5pt]{ESS: több BSS összekapcsolva, roaming lehetőséggel};
\end{tikzpicture}%
}
\caption{IEEE 802.11 infrastruktúrás hálózat: BSS, AP, DS és ESS}
\label{fig:zv23_wlan_bss_ess}
\end{figure}
